\paragraph{Supplement to A51: full deck transitivity in characteristic two.}
Let $C/k$ be a smooth projective geometrically connected genus-two
curve over a finite field of characteristic two. Let $D\to C$ be
a geometrically connected etale cover of odd degree $n>1$,
geometrically Galois with cyclic group generated by $\sigma$.
The deck group need not be constant. Write
\[
 F\sigma F^{-1}=\sigma^a,\qquad
 f=\operatorname{ord}_n(a)=\varphi(n),
 \tag{TD.1}
\]
where $F$ is arithmetic Frobenius on geometric points.
Thus $a$ acts transitively on all faithful characters.
Let $B=\operatorname{im}(ne_n:J_D\to J_D)$ be the actual
primitive conductor-$n$ subvariety. Then
\[
\begin{aligned}
 &f(B)=0\quad\hbox{or}\quad f(B)=\dim B=\varphi(n),\\
 &B(k)[2]=0\quad\Longleftrightarrow\quad f(B)=0,\\
 &f(B)=\varphi(n)\quad\Longrightarrow\quad
                         \dim_{\mathbf F_2}B(k)[2]=1.
\end{aligned}
 \tag{TD.2}
\]
Here $f(B)$ denotes geometric two-rank. In particular, in this
fully transitive case the rational primitive-unit test of A46
is equivalent to zero geometric primitive two-rank. Full
arithmetic transitivity does not remove the rational two-torsion
of an ordinary primitive factor.

\begin{proof}
Use the full cyclotomic algebra and its reduction
\[
 \mathcal O=\mathbf Z_2[Z]/\Phi_n(Z),\qquad
 A=\mathcal O/2\mathcal O,\qquad \tau(Z)=Z^a.
 \tag{TD.3}
\]
Since $n$ is odd, $A$ is a finite product of finite fields.
After extending to $\overline{\mathbf F}_2$ its factors are
indexed by all primitive $n$th roots. Condition (TD.1) makes
$\tau$ transitive on those indices, so
\[
              \dim_{\mathbf F_2}A=f,\qquad A^\tau=\mathbf F_2.
 \tag{TD.4}
\]
In particular it permutes the field factors of $A$ transitively.

The actual geometric group $V=B[2](\bar k)$ is an $A$-module
by (SD.8)--(SD.9a), and $F$ is a bijective
$\tau$-semilinear map on $V$.
Semilinearity makes the ranks of $V$ over the field factors
of $A$ equal. Consequently $V$ is free of some rank $r$ over
$A$ and $\dim_{\mathbf F_2}V=rf$.
The genus-two unramified character formula gives
$\dim B=\varphi(n)=f$: each nontrivial character has degree
two in the zeta numerator, and the actual primitive image
retains exactly the faithful characters
\cite[B.2]{A1run}; see also (OQ.7)--(OQ.8).
The two-rank bound $f(B)\le\dim B$
\cite[Remark~8.5, the bound $r\le g$]{MilneAV}
therefore forces $r=0$ or $r=1$.

If $r=0$, there is no geometric or rational two-torsion.
If $r=1$, choose an $A$-basis of $V$. In that basis
$F=c\tau$ for a unit $c\in A^\times$, and hence
\[
 F^f=c\tau(c)\cdots\tau^{f-1}(c)
                         \in(A^\tau)^\times=\mathbf F_2^\times=
                         \{1\}.
 \tag{TD.5}
\]
The finite etale torsor descent already proved in A45,
(SD.3)--(SD.6), now applies to this module killed by two.
It gives
\[
                  V\simeq A\otimes_{\mathbf F_2}V^F.
 \tag{TD.6}
\]
Equivalently, this is the norm-one Hilbert90 calculation for
the full finite etale product algebra. No hypothesis that
$A$ is a field, and no division by $f$, is needed.
Comparing dimensions in (TD.6) gives
$\dim_{\mathbf F_2}V^F=1$.
Since $V^F=B(k)[2]$, all assertions of (TD.2) follow.
The descent input is the same unrestricted module theorem
\cite[Proposition~35.3.9, Tag~\texttt{023N}]{Stacks023N}.
\end{proof}

The determinant makes the obstruction equally explicit.
Put $q=\#k$. Over an auxiliary-adic splitting field the
faithful summands of the Tate module have dimension two,
and $F$ cyclically permutes all $f$ of them.
Transport a basis successively by $F$ around this cycle;
the successive block maps are identities except for the last,
which is the return matrix $R$ on the first summand.
Eliminating these successive blocks in $1-TF$ gives
$\det(1-TF)=\det(1-T^fR)$. Therefore the
actual zeta numerator of $B$ has the form
\[
                       H_B(T)=1+S T^f+q^fT^{2f},\qquad S\in\mathbf Z.
 \tag{TD.7}
\]
The integrality and top coefficient are the Frobenius
polynomial identities for the $f$-dimensional abelian variety
$B$ \cite[Theorem~19.1]{MilneAV}.
In the ordinary case, (TD.2) makes
$\#B(k)=H_B(1)=1+S+q^f$ even, so $S$ is odd and
\[
                        \overline{H_B(T)}=1+T^f.
 \tag{TD.8}
\]
Thus the reduced return scalar equals one, rather than being
forced away from one by polarization or canonical symmetry.
The rational two-primary group has one generator and its
exponent is exactly
$2^{v_2(1+S+q^f)}$; no bound on that positive valuation has
been obtained. In the zero-two-rank case the rational group
is zero and $S$ is even. These are statements about the
actual primitive factor, not abstract module realizations.

\paragraph{Why the unrestricted quotient supply does not target this tower.}
In the rank-one $\ell$-tower of (NC.10)--(NC.11), with
$n_r=\ell^r$ and minimal deck field $K_r=k_{f_r}$,
one has $f_r=t\ell^{r-s}$ for large $r$.
Over $K_r$, the prescribed pointed cover has a constant
cyclic Lang quotient of size $n_r$.
Set $Q_r=q^{f_r}$. The genus-two Weil bounds give
\cite[Theorem~11.1]{MilneJV}
\[
\begin{aligned}
 \#C(K_r)&\ge Q_r+1-4\sqrt{Q_r}>n_r
                                      &&\hbox{for all large }r,\\
 \#\ker\bigl(J_C(K_r)\twoheadrightarrow C_{n_r}\bigr)
   &\ge\frac{(\sqrt{Q_r}-1)^4}{n_r}\longrightarrow\infty.
\end{aligned}
 \tag{TD.9}
\]
The kernel bound uses the Jacobian Weil bound
\cite[Theorem~19.1]{MilneAV}.
Thus the empty-coset inequality $|H|>\#C(K_r)$ and the fixed
small kernel used in A49's unrestricted-quotient supply proof
are unavailable for this prescribed tower quotient.
That proof may still supply a unit character on another,
much larger quotient, with an uncontrolled conductor.

Norms from a common refinement constrain the inflated
characters and kill its faithful component, as (NC.8)--(NC.9)
show. Moreover (NC.15) bounds generators of each individual
primitive factor, not the sum of the new factors on all
levels or independent covers. Neither statement puts the
supplied unit orbit in the remaining component of the same
compatible two-point class. In the fully transitive ordinary
case, (TD.2)--(TD.8) specifically show that the presumed
missing rational two-torsion is present. The unresolved
input remains a useful zero-rank cover, a controlled actual
return valuation, or an additional pointed geometric
argument. No universal existence conclusion follows here.
